\documentclass{beamer}

\usepackage[T2A]{fontenc}
\usepackage[utf8]{inputenc}
\usepackage[english,russian]{babel}
\input{settings.tex}


\title{Введение}
\subtitle{Вычислительный интеллект, лекция 1}
\author{Сергей Савин}
\centering
\date{\mydate}



\begin{document}
\maketitle


%\begin{frame}{Content}
%
%\begin{itemize}
%\item Motivation
%\item Motivating example
%\begin{itemize}
%\item fmincon
%\item quadprog
%\item SVD-based solution
%\item CVX-based solution
%\end{itemize}
%\item Homework
%\end{itemize}
%
%\end{frame}

\begin{frame}{Мотивация}
	\begin{flushleft}
		
		Многие современные методы (например, планирование траекторий, оценка состояния, синтез управления) опираются на инструменты численной оптимизации. В этом курсе мы изучим численную оптимизацию с акцентом на выпуклые методы.
		
		\begin{block}{Чего мы хотим?}
			Перейти от принципа "надеюсь, это сработает" к глубокому пониманию математики и сценариев использования этих инструментов.
		\end{block}
		
		\begin{block}{Зачем нам это?}
			Это должно позволить нам решать гораздо более широкий круг задач, причем более эффективно.
		\end{block}
		
	\end{flushleft}
\end{frame}

\begin{frame}{Мотивирующий пример}
	\begin{flushleft}
		
		У нас есть следующая задача: найти такой $\mathbf x$, который минимизирует $\mathbf x^\top \mathbf M \mathbf x$, при условии $\mathbf C \mathbf x = \mathbf y$. Другими словами:
		
		\begin{equation} \label{eq:Example1:1}
			\begin{aligned}
				& \underset{\mathbf x}{\text{minimize}}
				& & \mathbf x^\top \mathbf M \mathbf x, \\
				& \text{subject to:}
				& & \mathbf C \mathbf x = \mathbf y.
			\end{aligned}
		\end{equation}
		
		Более конкретно:
		
		\begin{equation} \label{eq:Example1:2}
			\begin{aligned}
				& \underset{\mathbf x}{\text{minimize}}
				& & \begin{bmatrix} x_1 & x_2 & x_3 \end{bmatrix}
				\begin{bmatrix}
					1 & 0 & 1 \\
					0 & 5 & 0 \\
					1 & 0 & 3 \\
				\end{bmatrix}
				\begin{bmatrix} x_1 \\ x_2 \\ x_3 \end{bmatrix}, \\
				& \text{subject to:}
				& & \begin{bmatrix} 1 & 7 & 2 \end{bmatrix}
				\begin{bmatrix} x_1 \\ x_2 \\ x_3 \end{bmatrix}
				= 1.
			\end{aligned}
		\end{equation}
		
		Как нам ее решить?
		
	\end{flushleft}
\end{frame}

\begin{frame}[fragile]{Мотивирующий пример}
	\framesubtitle{fmincon}
	\begin{flushleft}
		
		Очень популярное решение — использовать солвер локальной оптимизации общего назначения, например \texttt{fmincon} из MATLAB. Вот одно из возможных решений:
		
		\begin{lstlisting}[language=Matlab]
			M = [1 0 1; 0 5 0; 1 0 3];
			C = [1 7 2];
			y = 1;
			fnc = @(x) x'Mx;
			con = @(x) deal([], C*x-y);
			x = fmincon(fnc, zeros(3, 1), [],[],[],[],[],[], con)
		\end{lstlisting}
		%
		Среднее время решения составляет \textbf{2.5} мс (это зависит от многих факторов, поэтому рассматривайте это только как относительную информацию).
		Решение: $\mathbf{x} = \begin{bmatrix} 0.0442 & 0.1239 & 0.0442 \end{bmatrix}$.
		
	\end{flushleft}
\end{frame}


\begin{frame}[fragile]{Мотивирующий пример}
	\framesubtitle{quadprog}
	\begin{flushleft}
		
		Более изощренный, но все еще очень прямолинейный подход — использовать специализированный решатель для данного класса задач — \texttt{quadprog} из MATLAB. Вот решение:
		
		\begin{lstlisting}[language=Matlab]
			M = [1 0 1; 0 5 0; 1 0 3];
			C = [1 7 2];
			y = 1;
			
			x = quadprog(2*M,[],[],[],C,y)
		\end{lstlisting}
		
		Среднее время решения составляет \textbf{0.3} мс, что на порядок меньше, чем с \texttt{fmincon}.
		
	\end{flushleft}
\end{frame}

\begin{frame}[fragile]{Мотивирующий пример}
	\framesubtitle{Решение на основе SVD}
	\begin{flushleft}
		
		Мы можем использовать алгебраическое решение, основанное на SVD-разложении (через ядро и псевдообратную матрицу), следующим образом:
		
		\begin{lstlisting}[language=Matlab]
			M = [1 0 1; 0 5 0; 1 0 3];
			C = [1 7 2];
			y = 1;
			
			tol = 10^(-5);
			[P, N] = pinv_null(C, tol);
			x = ( eye(3) - N((N'MN) \ (N'M)) ) * P*y
		\end{lstlisting}
		
		Где \texttt{pinv\_null} — это функция, объединяющая \texttt{pinv} и \texttt{null}, полученная из одного SVD-разложения.
		
		Среднее время решения составляет \textbf{0.016} мс, что $\sim$ в 20 раз быстрее, чем \texttt{quadprog}, и $\sim$ в 150 раз быстрее, чем \texttt{fmincon}.
		
	\end{flushleft}
\end{frame}

\begin{frame}[fragile]{Мотивирующий пример}
	\framesubtitle{Левое матричное деление}
	\begin{flushleft}
		
		Мы можем использовать специализированный инструмент для решения линейных уравнений — \texttt{mldivide}, который в нашем случае применит LU-разложение (размер матрицы меньше 16x16). Вот решение:
		
		\begin{lstlisting}[language=Matlab]
			T1 = M \ C';
			t2 = (C*T1) \ y;
			x = T1*t2;
		\end{lstlisting}
		
		Среднее время решения составляет \textbf{0.0065} мс, что в \textbf{380} раз быстрее, чем \texttt{fmincon}, в \textbf{45} раз быстрее, чем \texttt{quadprog}, и в \textbf{2.5} раза быстрее, чем SVD.
		
	\end{flushleft}
\end{frame}

\begin{frame}[fragile]{Мотивирующий пример}
	\framesubtitle{Решение на основе CVX}
	\begin{flushleft}
		
		{\small 
		
		Наконец, мы можем задействовать один из самых мощных инструментов выпуклой оптимизации с удобным для пользователя стилем кодирования — CVX:
		
		\begin{lstlisting}[language=Matlab]
			M = [1 0 1; 0 5 0; 1 0 3];
			C = [1 7 2];
			y = 1;
			cvx_begin
			variables x(3);
			minimize( x' * M * x );
			subject to
			C*x == y;
			cvx_end
		\end{lstlisting}
		
		Однако мы увидим, что накладные расходы на вызов решателя для этой задачи чрезмерны. Среднее время решения составляет \textbf{219} мс, что $\sim$ в 90 раз медленнее, чем \texttt{fmincon}.
		
	}
	\end{flushleft}
\end{frame}




%\begin{frame}{Homework}
%% \framesubtitle{Parameter estimation}
%\begin{flushleft}
%
%Solve problem \eqref{eq:Example1:2} by the method you know and understand.
%
%\end{flushleft}
%\end{frame}

\myqrframe

\end{document}
